\begin{frontmatter}              % The preamble begins here.


%\pretitle{Pretitle}
\title{Chemical Identity Lost in Regulation\@:\\A Study of Semantic Interoperability in European Chemical Substance Data}
\runningtitle{Chemical Identity Lost in Regulation}
%\subtitle{Subtitle}

\ifblindreview\else
\author[A]{\fnms{Geert} \snm{Van Haute}\orcid{0009-0002-9836-8139}%
\thanks{Corresponding Author: Geert Van Haute, contact details.}},
\author[A]{\fnms{Pieter} \snm{Fannes}\orcid{....-....-....-....}}
,
\author[A]{\fnms{Stijn} \snm{Goedertier}\orcid{....-....-....-....}}
and
\author[A]{\fnms{Maxim} \snm{Van de Wynckel}\orcid{0000-0003-0314-7107}}

\runningauthor{G. Van Haute et al.}
\address[A]{Department of Environment and Spatial Development, Flemish Government, Belgium}
\fi


\begin{abstract}
    Chemical substance identity is inconsistently represented across European regulatory datasets, creating fundamental barriers to semantic interoperability. While structure-based identifiers such as InChI and InChIKey unambiguously define chemical entities, regulatory systems rely on heterogeneous identifiers including CAS numbers, EC numbers, textual names, and group descriptions that may represent mixtures, substance groups, or analytical parameters rather than molecules with a uniquely defined chemical structure.

    In this research, we analysed interoperability challenges across 18 European regulatory datasets and demonstrate that; (i)~37.5\% of entries cannot be linked to a defined chemical structure, (ii)~CAS numbers exhibit both one-to-many and many-to-one mappings with molecular structures, and (iii)~embedding-based classification fails for non-structure substances, with the best model achieving a ChemOnt Hit@1 of only 5.75\%. The combinatorial complexity of cross-list identifier mappings makes retrospective harmonisation infeasible at scale.

    These findings demonstrate that interoperability cannot be achieved through post hoc technical solutions alone. We argue that it must be addressed at the level of data modelling and legislation, by adopting structure-based identifiers, formal ontologies, and machine-readable representations directly in regulatory data standards.
\end{abstract}



\begin{keyword}
chemical substance identification\sep semantic interoperability\sep regulatory interoperability\sep REACH\sep
InChIKey\sep cheminformatics\sep FAIR data\sep knowledge graph\sep
European environmental legislation\sep Interoperable Europe
\end{keyword}
\end{frontmatter}
